% Bagian: normal-modal-logic
% Bab: tableaux
% Subbagian: rules-for-K

\documentclass[../../../include/open-logic-section]{subfiles}

\begin{document}

\olfileid[id]{nml}{tab}{rul}

\olsection{Aturan bagi \Ax{K}}

Aturan bagi konektif proposisional biasa sama dengan aturan bagi
!!{tableau} proposisional bertanda biasa, hanya dengan prefiks yang
ditambahkan. Dalam setiap kasus, aturan yang diterapkan pada !!{formula}
bertanda $\sFmla{S}{!A}[\sigma]$ menghasilkan !!{formula} baru yang juga
berprefiks~$\sigma$. Hal ini seharusnya jelas secara intuitif: misalnya, jika
$!A \land !B$ benar pada (dunia yang dinamai oleh)~$\sigma$, maka $!A$ dan
$!B$ benar pada~$\sigma$ (dan bukan pada dunia lain). Kita menghimpun aturan
proposisional dalam \olref{tab:prop-rules}.

\begin{table}
  \[\def\arraystretch{3}\begin{array}{|c|c|}
    \hline
    \AxiomC{\sFmla{\True}{\lnot !A}[\sigma]}
    \RightLabel{\TRule{\True}{\lnot}}
    \UnaryInfC{\sFmla{\False}{!A}[\sigma]}
    \DisplayProof
    &
    \AxiomC{\sFmla{\False}{\lnot !A}[\sigma]}
    \RightLabel{\TRule{\False}{\lnot}}
    \UnaryInfC{\sFmla{\True}{!A}[\sigma]}
    \DisplayProof
    \\[1ex]
    \hline
    \AxiomC{\sFmla{\True}{!A \land !B}[\sigma]}
    \RightLabel{\TRule{\True}{\land}}
    \UnaryInfC{\sFmla{\True}{!A}[\sigma]}
    \noLine
    \UnaryInfC{\sFmla{\True}{!B}[\sigma]}
    \DisplayProof
    &
    \AxiomC{\sFmla{\False}{!A \land !B}[\sigma]}
    \RightLabel{\TRule{\False}{\land}}
    \UnaryInfC{$\sFmla{\False}{!A}[\sigma] \quad \mid \quad
      \sFmla{\False}{!B}[\sigma]$}
    \DisplayProof
    \\[2ex]
    \hline
    \AxiomC{\sFmla{\True}{!A \lor !B}[\sigma]}
    \RightLabel{\TRule{\True}{\lor}}
    \UnaryInfC{$\sFmla{\True}{!A}[\sigma] \quad \mid \quad
      \sFmla{\True}{!B}[\sigma]$}
    \DisplayProof
    &
    \AxiomC{\sFmla{\False}{!A \lor !B}[\sigma]}
    \RightLabel{\TRule{\False}{\lor}}
    \UnaryInfC{\sFmla{\False}{!A}[\sigma]}
    \noLine
    \UnaryInfC{\sFmla{\False}{!B}[\sigma]}
    \DisplayProof
    \\[2ex]
    \hline
    \AxiomC{\sFmla{\True}{!A \lif !B}[\sigma]}
    \RightLabel{\TRule{\True}{\lif}}
    \UnaryInfC{$\sFmla{\False}{!A}[\sigma] \quad \mid
      \quad \sFmla{\True}{!B}[\sigma]$}
    \DisplayProof
    &
    \AxiomC{\sFmla{\False}{!A \lif !B}[\sigma]}
    \RightLabel{\TRule{\False}{\lif}}
    \UnaryInfC{\sFmla{\True}{!A}[\sigma]}
    \noLine
    \UnaryInfC{\sFmla{\False}{!B}[\sigma]}
    \DisplayProof
    \\[2ex]
    \hline
  \end{array}\]
  \caption{Aturan !!{tableau} berprefiks bagi konektif proposisional}
  \ollabel{tab:prop-rules}
\end{table}

Syarat penutupan sama seperti syarat bagi !!{tableau} biasa, meskipun kita
mensyaratkan agar bukan hanya !!{formula}, melainkan prefiksnya juga cocok.
Jadi, suatu cabang tertutup jika memuat keduanya,
\[
\sFmla{\True}{!A}[\sigma] \quad\text{dan}\quad
\sFmla{\False}{!A}[\sigma]
\]
untuk suatu prefiks $\sigma$ dan !!{formula}~$!A$.

Aturan untuk menetapkan asumsi juga sama seperti aturan bagi !!{tableau}
biasa, kecuali bahwa untuk asumsi kita selalu menggunakan prefiks~$1$.
(Tidak penting prefiks mana yang kita gunakan, asalkan prefiks tersebut sama
untuk semua asumsi.) Jadi, misalnya, kita mengatakan bahwa
\[
!B_1, \dots, !B_n \Proves !A
\]
jika dan hanya jika terdapat tableau tertutup untuk asumsi-asumsi
\[
\sFmla{\True}{!B_1}[1], \dots, \sFmla{\True}{!B_n}[1],
\sFmla{\False}{!A}[1].
\]

Untuk operator modal\iftag{prvBox}{\iftag{prvDiamond}{~$\Box$
    dan}{~$\Box$}}{}\iftag{prvDiamond}{~$\Diamond$}{}, prefiks pada
kesimpulan aturan yang diterapkan pada !!a{formula} berprefiks~$\sigma$
adalah $\sigma.n$. Namun, nilai $n$ yang diizinkan bergantung pada apakah
tandanya~$\True$ atau~$\False$.

\iftag{prvBox}{Aturan $\TRule{\True}{\Box}$ memperluas suatu cabang
  yang memuat $\sFmla{\True}{\Box !A}[\sigma]$ dengan
  $\sFmla{\True}{!A}[\sigma.n]$.\iftag{prvDiamond}{ Demikian pula,
    a}{}}{A}\iftag{prvDiamond}{turan $\TRule{\False}{\Diamond}$
  memperluas suatu cabang yang memuat
  $\sFmla{\False}{\Diamond !A}[\sigma]$ dengan
  $\sFmla{\False}{!A}[\sigma.n]$.}{}
\iftag{notprvBox,notprvDiamond}{Aturan itu}{Aturan-aturan itu} hanya dapat
diterapkan untuk prefiks~$\sigma.n$ yang \emph{sudah} muncul pada cabang
tempat aturan tersebut diterapkan. Mari kita sebut prefiks semacam itu
``sudah digunakan'' (pada cabang tersebut).

\iftag{prvBox}{Aturan $\TRule{\False}{\Box}$ memperluas suatu cabang
  yang memuat $\sFmla{\False}{\Box !A}[\sigma]$ dengan
  $\sFmla{\False}{!A}[\sigma.n]$.\iftag{prvDiamond}{ Demikian pula,
    a}{}}{A}\iftag{prvDiamond}{turan $\TRule{\True}{\Diamond}$ memperluas
  suatu cabang yang memuat $\sFmla{\True}{\Diamond !A}[\sigma]$ dengan
  $\sFmla{\True}{!A}[\sigma.n]$.}{}
Namun, \iftag{notprvBox,notprvDiamond}{aturan ini}{aturan-aturan ini} hanya
dapat diterapkan untuk prefiks~$\sigma.n$ yang \emph{belum} muncul pada
cabang tempat aturan tersebut diterapkan. Kita menyebut prefiks semacam itu
``baru'' (bagi cabang tersebut).

Aturan-aturan itu diberikan dalam \olref{tab:rules-K}.

\begin{table}
  \begin{center}
    \def\arraystretch{3}\def\fCenter{}
    \begin{tabular}{|c|c|}
    \hline
    \iftag{prvBox}{
      \AxiomC{\sFmla{\True}{\Box !A}[\sigma]}
      \RightLabel{\TRule{\True}{\Box}}
      \UnaryInfC{\sFmla{\True}{!A}[\sigma.n]}
      \DisplayProof
      &
      \AxiomC{\sFmla{\False}{\Box !A}[\sigma]}
      \RightLabel{\TRule{\False}{\Box}}
      \UnaryInfC{\sFmla{\False}{!A}[\sigma.n]}
      \DisplayProof\\
      $\sigma.n$ sudah digunakan & $\sigma.n$ baru
      \\[1ex]
      \hline}{}
    \iftag{prvDiamond}{
      \AxiomC{\sFmla{\True}{\Diamond !A}[\sigma]}
      \RightLabel{\TRule{\True}{\Diamond}}
      \UnaryInfC{\sFmla{\True}{!A}[\sigma.n]}
      \DisplayProof
      &
      \AxiomC{\sFmla{\False}{\Diamond !A}[\sigma]}
      \RightLabel{\TRule{\False}{\Diamond}}
      \UnaryInfC{\sFmla{\False}{!A}[\sigma.n]}
      \DisplayProof\\
      $\sigma.n$ baru & $\sigma.n$ sudah digunakan
      \\[1ex]
      \hline}{}
    \end{tabular}
  \end{center}
  \caption{Aturan modal bagi \Ax{K}.}
  \ollabel{tab:rules-K}
\end{table}

Syarat bahwa prefiks bagi
\iftag{prvBox}{\TRule{\True}{\Box}}{\TRule{\False}{\Diamond}} harus sudah
digunakan memang diperlukan; tanpa syarat itu, kita akan menganggap yang
berikut sebagai !!{tableau} tertutup:
\iftag{notprvBox,notprvDiamond}{%
  \iftag{prvBox}{%
  \begin{oltableau}
    [\pFmla{\True}{\Box \formula{A}}{1}, just = \TAss
      [\pFmla{\False}{\lnot\Box\lnot \formula{A}}{1}, just = \TAss
        [\pFmla{\True}{\formula{A}}{1.1}, just = {\TRule{\True}{\Box}[1]}
          [\pFmla{\True}{\Box\lnot\formula{A}}{1},
            just ={\TRule{\False}{\lnot}[2]}
            [\pFmla{\True}{\lnot\formula{A}}{1.1},
              just = {\TRule{\True}{\Box}[4]}
              [\pFmla{\False}{\formula{A}}{1.1},
                  just ={\TRule{\True}{\lnot}[5]}, close]
            ]
          ]
        ]
      ]
    ]
  \end{oltableau}
  }{
  \begin{oltableau}
    [\pFmla{\True}{\lnot\Diamond\lnot \formula{A}}{1}, just = \TAss
      [\pFmla{\False}{\Diamond\formula{A}}{1}, just = \TAss
        [\pFmla{\False}{\formula{A}}{1.1}, just = {\TRule{\False}{\Diamond}[2]}
          [\pFmla{\False}{\Diamond\lnot\formula{A}}{1},
            just ={\TRule{\True}{\lnot}[1]}
            [\pFmla{\False}{\lnot\formula{A}}{1.1},
              just = {\TRule{\False}{\Diamond}[4]}
              [\pFmla{\True}{\formula{A}}{1.1},
                  just ={\TRule{\True}{\lnot}[5]}, close]
            ]
          ]
        ]
      ]
    ]
  \end{oltableau}
}}{
  \begin{oltableau}
    [\pFmla{\True}{\Box \formula{A}}{1}, just = \TAss
      [\pFmla{\False}{\Diamond \formula{A}}{1}, just = \TAss
        [\pFmla{\True}{\formula{A}}{1.1}, just = {\TRule{\True}{\Box}[1]}
          [\pFmla{\False}{\formula{A}}{1.1},
            just = {\TRule{\False}{\Diamond}[2]}, close]
        ]
      ]
    ]
\end{oltableau}}

Namun, $\Box \formula{A} \Entails/ \Diamond \formula{A}$, sehingga sistem
pembuktian kita akan menjadi tidak sahih. Demikian pula,
$\Diamond \formula{A} \Entails/ \Box \formula{A}$; tetapi tanpa syarat
bahwa prefiks bagi
\iftag{prvBox}{\TRule{\False}{\Box}}{\TRule{\True}{\Diamond}} harus baru,
yang berikut akan menjadi tableau tertutup:
\iftag{notprvBox,notprvDiamond}{%
  \iftag{prvBox}{%
  \begin{oltableau}
    [\pFmla{\True}{\lnot\Box\lnot \formula{A}}{1}, just = \TAss
      [\pFmla{\False}{\Box \formula{A}}{1}, just = \TAss
        [\pFmla{\False}{\formula{A}}{1.1}, just = {\TRule{\True}{\Box}[2]}
          [\pFmla{\False}{\Box\lnot\formula{A}}{1},
            just ={\TRule{\True}{\lnot}[1]}
            [\pFmla{\False}{\lnot\formula{A}}{1.1},
              just = {\TRule{\False}{\Box}[4]}
              [\pFmla{\True}{\formula{A}}{1.1},
                  just ={\TRule{\False}{\lnot}[5]}, close]
            ]
          ]
        ]
      ]
    ]
  \end{oltableau}
  }{
  \begin{oltableau}
    [\pFmla{\True}{\Diamond \formula{A}}{1}, just = \TAss
      [\pFmla{\False}{\lnot\Diamond\lnot\formula{A}}{1}, just = \TAss
        [\pFmla{\True}{\formula{A}}{1.1}, just = {\TRule{\True}{\Diamond}[1]}
          [\pFmla{\True}{\Diamond\lnot\formula{A}}{1},
            just ={\TRule{\False}{\lnot}[2]}
            [\pFmla{\True}{\lnot\formula{A}}{1.1},
              just = {\TRule{\True}{\Diamond}[4]}
              [\pFmla{\False}{\formula{A}}{1.1},
                  just ={\TRule{\True}{\lnot}[5]}, close]
            ]
          ]
        ]
      ]
    ]
  \end{oltableau}
}}{
  \begin{oltableau}
    [\pFmla{\True}{\Diamond \formula{A}}{1}, just = \TAss,name=one
      [\pFmla{\False}{\Box \formula{A}}{1}, just = \TAss, name=two
        [\pFmla{\True}{\formula{A}}{1.1}, just = {\TRule{\True}{\Diamond}[1]}
          [\pFmla{\False}{\formula{A}}{1.1},
            just = {\TRule{\False}{\Box}[2]}, close]
        ]
      ]
    ]
  \end{oltableau}}

\end{document}
